% !TeX spellcheck = en_US
% This is a derived work from 
% http://en.wikibooks.org/wiki/LaTeX/Title_Creation#A_practical_example

\begin{titlepage}
\begin{center}

\includegraphics[width=0.15\textwidth]{img/unibz-logo.pdf}
~\\[1cm]
\textsc{\LARGE Free University of Bolzano/Bozen}
\\[1.5cm]
\textsc{\Large Report for the Course \\ Semantic Technologies}
\\[0.5cm]

% Title
\HRule{0.5pt} 
\\[0.4cm]
{ 
	\huge 
	\bfseries 
	Semantic Techniques to find new dimensions and attributes for a Tourist Inquiry Data Warehouse
	\\[0.4cm] 
}
\HRule{2pt} 
\\[1.5cm]

% Author and supervisor
\noindent
\begin{minipage}[t][][b]{0.4\textwidth}
	\begin{flushleft} \large
		\emph{Authors:}\\
		Peter \textsc{Moser}
	\end{flushleft}
\end{minipage}%
\begin{minipage}[t][][b]{0.4\textwidth}
	\begin{flushright} \large
		\emph{Submitted to:} \\
		Prof.~Enrico \textsc{Franconi}, Dr.~Guohui \textsc{Xiao}
	\end{flushright}
\end{minipage}


\vspace{.8cm}

{\large \today}

% Bottom of the page
\vfill

\HRule{0.1pt} 
\begin{flushleft}
This report shows the design, implementation, and queries of a tourist inquiries data warehouse through a semantic data model. In detail it contains, first, a domain analysis with defined business goals, a conceptual design with dimensions, facts, and a star schema, second, the architecture of the system with explored possibilities and failures, third, I explain the semantic techniques and ontology used. Finally, I present a Java application, with which an user can query the ontology in two ways: First, by writing a SPARQL query to retrieve data from a tourist inquiries ontology via the ONTOP framework, and second, to query a remote SPARQL endpoint through rdf4j.  
\end{flushleft}
\HRule{0.1pt}
\end{center}
\end{titlepage}